% Ziren as a processor block diagram: units (chips) and buses (interactions).
\begin{tikzpicture}[
  every node/.style={font=\small},
  unit/.style={rectangle, draw=zirenblue!85, rounded corners=2pt, text width=2.0cm, minimum height=0.7cm, align=center, fill=zirenblue!7, inner sep=2pt},
  accel/.style={unit, draw=zirenorange!90, fill=zirenorange!10},
  mem/.style={rectangle, rounded corners=2pt, draw=zirenteal!85, text width=2.7cm, minimum height=0.7cm, align=center, fill=zirenteal!8, inner sep=2pt},
  rom/.style={rectangle, rounded corners=2pt, draw=zirengold!90!black, text width=2.7cm, minimum height=0.7cm, align=center, fill=zirengold!12, inner sep=2pt},
  bus/.style={line width=1.15pt, zirenslate!85},
  programbus/.style={bus, zirenblue!90},
  statebus/.style={bus, zirenteal!90!black},
  memorybus/.style={bus, zirenorange!90!black},
  globalbus/.style={bus, zirenslate!95},
  buslab/.style={font=\small, black!70, right},
  tap/.style={-{Latex[length=1.6mm]}, thin, zirenslate!90},
]
  % instruction units (functional units), pitch 2.3 cm
  \node[unit] (alu)   at (0,0)    {ALU\\\texttt{AddSub}, \texttt{Bitwise}};
  \node[unit] (shift) at (2.3,0)  {Shift\\\texttt{ShiftLeft}\\\texttt{(Imm)}};
  \node[unit] (mul)   at (4.6,0)  {Mul/Div\\\texttt{Mul}, \texttt{DivRem}};
  \node[unit] (ctl)   at (6.9,0)  {Control\\\texttt{Branch}, \texttt{Jump}};
  \node[unit] (ld)    at (9.2,0)  {Load/Store\\\texttt{LoadWord}, \ldots};
  \node[unit] (sys)   at (11.5,0) {Syscall\\\texttt{Syscall}\\\texttt{Instrs}};
  \node[accel, text width=2.5cm] (pre) at (14.9,0) {Accelerators\\Keccak, SHA, EC, \ldots};

  % buses
  \draw[programbus] (-1.2,1.35) -- (16.4,1.35) node[buslab] {program bus};
  \draw[statebus] (-1.2,1.0)  -- (16.4,1.0)  node[buslab] {state bus};
  \draw[memorybus] (-1.2,-1.0) -- (16.4,-1.0) node[buslab] {memory bus};
  \draw[bus] (-1.2,-1.35) -- (16.4,-1.35) node[buslab] {byte / range bus};
  \draw[globalbus] (-1.2,-1.7) -- (16.4,-1.7) node[buslab] {global bus};
  \foreach \u in {alu,shift,mul,ctl,ld,sys,pre} {
    \draw[tap] (\u.north) -- ++(0,0.5);
    \draw[tap] (\u.south) -- ++(0,-0.5);
  }
  \draw[tap] (sys.east) -- node[above=1pt, font=\scriptsize, black!70, align=center] {syscall\\bus} (pre.west);

  % tables and memory units
  \node[rom] (prog) at (1.2,2.5)  {\texttt{Program} ROM\\(preprocessed, in vk)};
  \node[rom, text width=3.4cm] (pv)   at (5.6,2.5)  {Public values\\entry/exit state, digests};
  \node[mem] (bump) at (0.6,-2.9)  {\texttt{MemoryBump}\\register anchors};
  \node[rom] (byte) at (4.0,-2.9)  {\texttt{Byte} / \texttt{Range} tables\\(preprocessed)};
  \node[mem] (loc)  at (7.6,-2.9)  {\texttt{MemoryLocal}\\shard first/last touch};
  \node[mem] (glob) at (11.0,-2.9) {\texttt{Global}\\septic digest};
  \node[mem] (init) at (14.6,-2.9) {\texttt{MemoryGlobal}\\\texttt{Init}/\texttt{Finalize}};

  \draw[tap] (prog.south) -- (1.2,1.35);
  \draw[tap] (pv.south) -- (5.6,1.0);
  \draw[tap] (bump.north) -- (0.6,-1.0);
  \draw[tap] (byte.north) -- (4.0,-1.35);
  \draw[tap] (loc.north) -- (7.6,-1.0);
  \draw[tap] (loc.east) -- (glob.west);
  \draw[tap] (glob.north) -- (11.0,-1.7);
  \draw[tap] (init.north) -- (14.6,-1.7);
\end{tikzpicture}
